\documentclass[11pt]{article}

\usepackage{fontspec}
\setmainfont{Times New Roman}
\usepackage[margin=1in]{geometry}
\usepackage{amsmath}
\usepackage{booktabs}
\usepackage{hyperref}
\usepackage{url}
\usepackage{microtype}
\sloppy

\hypersetup{
  colorlinks=true,
  linkcolor=black,
  urlcolor=blue,
  citecolor=black
}

\title{Детерминированный аудит утверждений, написанных языковой моделью}
\author{Владимир В.~Боровиков\\
\small Независимый исследователь\\
\small \url{https://orcid.org/0009-0005-7773-3581}}
\date{Октябрь 2026}

\begin{document}
\maketitle

\begin{abstract}
Языковые модели всё чаще пишут описательную часть исследований: отчёты,
вердикты, разделы обсуждения. Опасен при этом не очевидный вымысел, а
\emph{правдоподобное преувеличение}: число верно, интерпретация нет, а текст
звучит уверенно. Мы описываем протокол, который делает такое преувеличение
механически обнаружимым, и сообщаем, что получилось, когда мы применили его к
собственному корпусу из 1\,416 проб в десяти экспериментальных сериях. Аудит
покрывает десять $p$-значений и двадцать одну дробь; его границы оговорены в
разделе~\ref{sec:status}.

Протокол опирается на шесть требований: заморозка до первой пробы, внешняя
память вместо припоминания из контекста, разделение генератора и аудитора,
детерминированная сверка каждого приведённого числа, трёхуровневая таксономия
ошибок утверждений и состязательный обзор несколькими моделями. Мы предлагаем
исполняемый чеклист из двенадцати пунктов и честно сообщаем его статус на
собственной работе: шесть пройдено, один исправлен по процедуре, один не
пройден, два частично, два требуют человека. Протокол применён не полностью,
и чеклист это показывает.
\end{abstract}

\section{Введение}

Когда языковая модель описывает эксперимент, текст получается гладким по
построению. Гладкость не является свидетельством правильности, но читается как
свидетельство — и рецензентом, и автором. В наших сериях мы наблюдали
повторяющийся сюжет: первая редакция вердикта утверждала больше, чем позволял
архив, и отзывалась только после отдельного пересчёта.

Примеры из наших логов, все проверяемы по открытому репозиторию:

\begin{itemize}
  \item нулевой результат ($p = 0{,}338$) описан как опровержение «со 100\%
        строгостью»;
  \item сравнение $2/72$ против $0/72$ названо «законом» (точный критерий
        Фишера $p = 0{,}4965$~\cite{fisher1935});
  \item $p$-значение напечатано как $4{,}81\times10^{-13}$, тогда как точный
        двусторонний критерий Фишера даёт $2{,}909\times10^{-18}$, — ошибка
        инструмента, а не модели, прожившая шесть недель;
  \item сноска отправляла читателя в файл данных, где не было ни одного из
        цитируемых байтов.
\end{itemize}

Последние два случая показательны. Третий породил наш собственный
проверяльщик; четвёртый — модель, уверенно написавшая фразу о файле, который
не открывала. Ни то, ни другое не поймано вычиткой.

\section{Таксономия ошибок утверждений}

Мы различаем три яруса, потому что ловятся они по-разному.

\textbf{Ярус 0 — фабрикация референта.} Текст ссылается на объект, которого
нет: файл, строку, набор данных, значение. Числа при этом могут быть
арифметически безупречны. Ни пересчёт, ни вычитка этого не обнаруживают,
потому что пересчитывать нечего. Нужна проверка существования.

\textbf{Ярус I — интерпретационный перебор.} Архив подтверждает число, но не
фразу. «$p = 0{,}34$» превращается в «опровергнуто». Лечится переформулировкой.

\textbf{Ярус II — числовая ошибка.} В архиве другое число, чем в тексте.
Лечится только пересчётом из лога.

Ярус 0 опаснее всех и защищён слабее всех: у первых двух есть механические
проверки, а нулевой исторически держался на одном внимании; протокол добавляет
для него механическую проверку существования (П5).

\section{Протокол}

\paragraph{П1. Заморозка.} Код, промпты, гипотезы и эндпоинты хэшируются до
первой пробы. Самоподписанный манифест — это декларация, а не заморозка: без
доверенной метки времени он доказывает лишь текущее состояние файлов. Мы
вводим статус \textsc{freeze-amended} для серий, чьи файлы изменились после
заморозки, но изменение задокументировано, результат пересчитан и совпал, а
регрессионные тесты пройдены. От «аннулировано» он отличается тем, что числа
сохраняются, от «заморожено» — тем, что побайтовое воспроизведение пайплайна
утрачено.

\paragraph{П2. Внешняя память.} Числа читаются из файлов и никогда не
припоминаются из контекста. Модель, которая должна помнить, рано или поздно
вспомнит неверно.

\paragraph{П3. Разделение контуров.} Генератор пишет, аудитор считает. Аудитор
независим только тогда, когда написан по отдельной спецификации; «другой
скрипт» — это ещё не независимость.

\paragraph{П4. Сверка.} Каждое приведённое число — $p$-значения, дроби,
проценты, границы интервалов — пересчитывается скриптом, и вывод скрипта
публикуется вместе со статьёй. Два дополнения оказались необходимы на
практике. Первое: \emph{тест на заведомо известный ответ} — намеренно
испорченный эталон, который проверяльщик обязан отвергнуть, иначе слепой
проверяльщик пропускает всё. Второе: статистику нельзя писать с нуля без
сверки; проверка, в которой одна самописная реализация сверяется с другой
самописной, сравнивает ошибку с ошибкой. Наша реализация критерия Фишера несла
абсолютный допуск $10^{-12}$, который при $p \approx 10^{-18}$ в
$3{,}4\times10^{5}$ раз больше самого значения и раздувал хвост на пять
порядков. Нашла её независимая реализация, а не протокол, — поэтому сверка
теперь обязательный шаг, а не счастливая случайность.

\paragraph{П5. Таксономия.} Как в разделе 2, причём Ярус~0 проверяется
скриптом, который извлекает из документа пути к файлам и диапазоны строк и
проверяет, что цели существуют.

\paragraph{П6. Обзор несколькими моделями.} Разные семейства моделей проверяют
друг друга. Это самый слабый столп: число проходов — рабочая договорённость, а
не доказанный оптимум, а модели, обученные на пересекающихся корпусах,
ошибаются коррелированно, поэтому «ансамбль» не гарантирует независимости
ошибок.

\section{Применение и честный статус}
\label{sec:status}

Чеклист состоит из двенадцати пунктов и исполним. На нашем корпусе он
показывает:

\begin{center}
\begin{tabular}{@{}lrl@{}}
\toprule
Статус & Кол-во & Что означает \\
\midrule
Пройдено   & 6 & проверено механически \\
Исправлено & 1 & правка после заморозки, задокументирована и пересчитана \\
Не пройден & 1 & независимость проверяльщика: писал тот же автор \\
Частично   & 2 & пре-регистрация, версионирование вердиктов \\
Вручную    & 2 & доверенная метка времени, неизменяемость архива \\
\bottomrule
\end{tabular}
\end{center}

Единственный непройденный пункт кодом не исправить. Нужен второй автор, который
напишет проверяльщик по спецификации, составленной не им. Мы говорим об этом
прямо, а не замазываем: чеклист, который можно пройти только соврав, хуже
отсутствия чеклиста.

Покрытие П4 — десять $p$-значений (девять двухвыборочных контрастов и
post-hoc пятибиновый тест), двадцать одна дробь и все приведённые
проценты, включая границы Клоппера–Пирсона~\cite{clopper1934}; всё пересчитано двумя независимыми
скриптами, совпадение которых проверяется при каждом запуске.

\section{Ограничения}

Выборка аудита неслучайна: мы проверяли те серии, где аудит проводился. О
сериях, которые не проверялись, мы ничего не утверждаем.

Девять гипотез восстановлены по тексту статьи, а не заморожены до прогона.
Поправка на множественность (Холм)~\cite{holm1979} обсуждается проспективно; к завершённым
сериям она не применялась.

Провенанс недоказуем: без доверенной метки времени мы не можем показать, что
архив был в нынешнем состоянии \emph{до} прогона, а не приведён к нему после.
Доказано более узкое: числа воспроизводятся из архива в том виде, в каком он
лежит в репозитории сегодня.

\section{Доступность}

Код, пробы и исполняемый чеклист:
\url{https://github.com/skywhiteangels-cell/llm-token-architect}.
Сопроводительное исследование: \url{https://doi.org/10.5281/zenodo.23813780}.

\paragraph{Идентификация модели.} Все приведённые здесь серии проб выполнялись
на \texttt{flash\_lite} — общедоступном тире модели; тексты вердиктов написаны
другой моделью общего назначения, что отмечено в архиве проекта. Ничего не
скрыто за проприетарными правами или соглашением о неразглашении: архив и
скрипты проверки опубликованы под лицензией MIT, поэтому любое число можно
пересчитать, не зная, какая модель его произвела.

\begin{thebibliography}{9}

\bibitem{clopper1934}
C.~J. Clopper and E.~S. Pearson.
\newblock The use of confidence or fiducial limits illustrated in the case of
  the binomial.
\newblock \emph{Biometrika}, 26(4):404--413, 1934.

\bibitem{fisher1935}
R.~A. Fisher.
\newblock \emph{The Design of Experiments}.
\newblock Oliver \& Boyd, Edinburgh, 1935.

\bibitem{holm1979}
S.~Holm.
\newblock A simple sequentially rejective multiple test procedure.
\newblock \emph{Scandinavian Journal of Statistics}, 6(2):65--70, 1979.

\end{thebibliography}

\end{document}
